\documentclass[10pt,a4paper]{amsart}
\usepackage[margin=19mm]{geometry}
\usepackage{amsmath,amssymb,mathtools}
\usepackage[scr=rsfs,cal=euler]{mathalfa}
\usepackage{xfrac}
\usepackage[unicode,hidelinks,bookmarks=false]{hyperref}
\newcommand{\ie}{i.e.\ }
\newcommand{\cf}{cf.\ }
\newcommand{\resp}{respectively\ }
\newcommand{\Subsection}{\subsection}
\providecommand{\nobreakdash}{\nobreak\hbox{-}}
\setlength{\emergencystretch}{2em}
\multlinegap0pt
\usepackage{amssymb}
\usepackage{enumerate}
\usepackage[shortlabels]{enumitem}
\usepackage{bm}
\usepackage{xcolor}
\usepackage{dsfont}
\usepackage{tikz}
\usepackage{algorithm}\usepackage{algpseudocode}
\newtheorem{assumption}{Assumption}
\newtheorem{lemma}{Lemma}
\newtheorem{proposition}{Proposition}
\usepackage{cleveref}
\crefname{assumption}{Assumption}{Assumptions}
\crefname{lemma}{Lemma}{Lemmas}
\crefname{proposition}{Proposition}{Propositions}
\newcommand{\R}{\mathbb{R}}
\title{Statistical Inference for Adversarial Training: Central Limit Theorems via Optimal Transport}
\author{Jakwang Kim, Dohyun Kwon}
\date{Source: arXiv:2609.22240v1 (CC BY 4.0)}
\begin{document}
\maketitle

\noindent\emph{Source and licence.} Statistical Inference for Adversarial Training: Central Limit Theorems via Optimal Transport. Version arXiv:2609.22240v1 (CC BY 4.0), distributed by arXiv under the Creative Commons Attribution 4.0 International licence, http://creativecommons.org/licenses/by/4.0/, as recorded in the arXiv licence metadata for this exact version and shown on the abstract page https://arxiv.org/abs/2609.22240v1. Full licence notice: \texttt{licenses/arxiv-2609.22240-CC-BY-4.0.txt}.\par\smallskip

\noindent\textit{Editorial scope.} Counted unit: the article's proposition potentials (source0.tex lines 1560--1591). The article's complete original proof of it is delivered with it (source0.tex lines 1593--1692, one \texttt{proof} environment of 100 lines). No mathematics is written, repaired or re-derived: the statement and the proof are the article's own lines, in the article's own notation, and no retained line is substituted or re-flowed.  Retained uncounted as the source-local context the proof needs: lines 862--873; lines 1329--1333; lines 1463--1477.  Nothing was omitted from any retained span, so the package carries no omission record.  No retained line contains a citation, so the wrapper carries no bibliography and no bibliography entry.  No retained line contains a figure environment, an image-inclusion command or drawing code, so no image is delivered and the package declares no asset dependency.  Editorial formatting only, in addition: an AMS \texttt{amsart} preamble that restates the article's own declarations and macros used by the retained lines, namely the environments \texttt{assumption}, \texttt{lemma}, \texttt{proposition} on separate counters, together with the standard packages those lines need.  The retained environments print in this document as Assumption 1, Lemma 1, Proposition 1 in order of appearance, whereas the article prints the same items with the numbering of its own full document.  The complete native source of the article as distributed in the arXiv e-print for this exact version is bundled unchanged as \texttt{sources/211019/source0.tex}.
\begin{assumption}\label{assumption: strict convexity of cost function}
$c(x,x')=h(x-x')$, where $h:\mathbb{R}^p \rightarrow [0, \infty)$ with $h(0)=0$ is a symmetric and continuous function satisfying
\begin{enumerate}
	\item[(A1)] $h$ is strictly convex on $\mathbb{R}^p $,
	\item[(A2)] given a height $r\in \R^+$ and an angle $\theta \in (0,\pi) $, there exists some $M:=M(r, \theta)>0$ such that for all $|p |>M$, one can find a cone 
	\begin{align*}
		K(r, \theta, z,p):=\left\lbrace x\in \mathbb{R}^p  : | x-p|| z|\cos(\theta/2)\leq \left< z,x-p \right>\leq  r| z| \right\rbrace,
	\end{align*}
	with vertex at $p$ on which $h$ attains its maximum at $p$,
	\item[(A3)] $\lim_{|x | \rightarrow \infty}\frac{h(x)}{|x | }= \infty $.
\end{enumerate}    
\end{assumption}

\begin{equation}\label{eq:mot_decomposed_dual}
\begin{aligned}
    \inf_{g=(g_i)_{i \in \mathcal{Y}} \in \mathcal{G}} \sum_{i \in \mathcal{Y}} \int (1-g_i(x_i)) d\mu_i(x_i)    
\end{aligned}
\end{equation}

\begin{lemma}
\label{lem: uniform regularity of potentials}
Assume that $c$ satisfies
\Cref{assumption: strict convexity of cost function}.
Then there exists a constant $L_c<\infty$, depending only on $c$,
such that, for every Borel measurable function
$f:\mathbb R^p\to[0,1]$, its $c$-transform
\[
    f^c(x)
    :=
    \inf_{y\in\mathbb R^p}
    \{f(y)+c(x,y)\}
\]
is $L_c$-Lipschitz.
\end{lemma}

\begin{proposition}[Stability of optimal potentials]
\label{prop: stability of optimal potentials}
Assume that $c$ satisfies
\Cref{assumption: strict convexity of cost function}.
Let $\eta^n,\eta\in\mathcal P(\mathcal X\times\mathcal Y)$ satisfy
\[
    \eta^n\rightharpoonup\eta
    \quad\text{as }n\to\infty.
\]
For each $n$, let
\[
    g^n=(g_i^n)_{i\in\mathcal Y}
\]
be an optimal potential for \eqref{eq:mot_decomposed_dual} with input
$\eta^n$, chosen so that, for every $i\in\mathcal Y$,
\[
    g_i^n=(f_i^n)^c
\]
for some Borel measurable function
$f_i^n:\mathbb R^p\to[0,1]$. Then every subsequence of $\{g^n\}$ admits a further subsequence
converging locally uniformly on $\mathbb R^p$ to an optimal potential
for \eqref{eq:mot_decomposed_dual} with input $\eta$.

In particular, if the optimal potential for the problem with input
$\eta$ is unique, say $g=(g_i)_{i\in\mathcal Y}$, then
\[
    g_i^n\longrightarrow g_i
    \quad
    \text{locally uniformly on }\mathbb R^p
\]
for every $i\in\mathcal Y$.
\end{proposition}

\begin{proof}
By \Cref{lem: uniform regularity of potentials}, the functions
$\{g_i^n\}_n$ are uniformly bounded and equi-Lipschitz on
$\mathbb R^p$. Hence, by the Arzelà--Ascoli theorem and a diagonal
argument, every subsequence of $\{g^n\}$ admits a further subsequence,
still denoted by $\{g^n\}$, such that
\[
    g_i^n\longrightarrow\bar g_i
    \quad
    \text{locally uniformly on }\mathbb R^p
\]
for every $i\in\mathcal Y$. Moreover,
$0\leq\bar g_i\leq1$ and $\bar g_i$ is $L_c$-Lipschitz.

The dual constraints are preserved under this convergence. Indeed, for
every $A\in S_{\mathcal Y}$ and
$x_A\in(\mathbb R^p)^A$,
\[
    \sum_{i\in A}g_i^n(x_i)
    \leq
    1+c_A(x_A).
\]
Passing to the limit gives
\[
    \sum_{i\in A}\bar g_i(x_i)
    \leq
    1+c_A(x_A).
\]
Thus $\bar g=(\bar g_i)_{i\in\mathcal Y}$ is feasible.

We next prove that $\bar g$ is optimal. Since $\mathcal Y$ is finite
and discrete, the weak convergence $\eta^n\rightharpoonup\eta$
implies
\[
    \eta_i^n\rightharpoonup\eta_i
    \quad
    \text{for every }i\in\mathcal Y.
\]
Let
$\widetilde g=(\widetilde g_i)_{i\in\mathcal Y}$ be any feasible
potential. Since $\sum_i\eta_i^n(1)=1$, the optimality of $g^n$ gives
\[
    \sum_{i\in\mathcal Y}\eta_i^n(g_i^n)
    \geq
    \sum_{i\in\mathcal Y}\eta_i^n(\widetilde g_i).
\]

We claim that
\[
    \eta_i^n(g_i^n)
    \longrightarrow
    \eta_i(\bar g_i)
\]
for every $i\in\mathcal Y$. Since
$\eta_i^n\rightharpoonup\eta_i$, the family
$\{\eta_i^n\}_n$ is tight. Given $\varepsilon>0$, choose a compact set
$K\subset\mathbb R^p$ such that
\[
    \sup_n\eta_i^n(K^c)<\varepsilon.
\]
Then, using $0\leq g_i^n,\bar g_i\leq1$,
\begin{align*}
    \left|
        \eta_i^n(g_i^n)-\eta_i(\bar g_i)
    \right|
    &\leq
    \sup_{x\in K}|g_i^n(x)-\bar g_i(x)|
    +\eta_i^n(K^c) \\
    &\quad
    +
    \left|
        \eta_i^n(\bar g_i)-\eta_i(\bar g_i)
    \right|.
\end{align*}
The first term converges to zero by local uniform convergence, while
the last term converges to zero because
$\bar g_i\in C_b(\mathbb R^p)$. Since $\varepsilon>0$ is arbitrary,
the claim follows.

Since each $\widetilde g_i\in C_b(\mathbb R^p)$, weak convergence also
gives
\[
    \eta_i^n(\widetilde g_i)
    \longrightarrow
    \eta_i(\widetilde g_i).
\]
Passing to the limit in the optimality inequality, we obtain
\[
    \sum_{i\in\mathcal Y}\eta_i(\bar g_i)
    \geq
    \sum_{i\in\mathcal Y}\eta_i(\widetilde g_i).
\]
Since $\widetilde g$ was arbitrary, $\bar g$ is an optimal potential
for the problem with input $\eta$.

Finally, suppose that the optimal potential for the limiting problem is
unique, say $g$. Then every subsequence of $\{g^n\}$ admits a further
subsequence converging locally uniformly to $g$. This implies that
the full sequence converges locally uniformly to $g$.
\end{proof}

\end{document}
